\chapterbackmatter{Solutions to Supplementary Exercises}

\begin{enumerate}
  \SupplementarySolution{1}{Discrete Symmetries and Dangerous MSSM Operators}
  \begin{enumerate}
    \item \emph{Matter-Parity Test of the Renormalizable Superpotential.}
    Gauge invariance and holomorphy allow the matter-parity-even terms
      \[
        W_{\rm MSSM}
          =\mu H_uH_d
           +(y_u)_{ij}Q_iH_u\bar U_j
           +(y_d)_{ij}Q_iH_d\bar D_j
           +(y_e)_{ij}L_iH_d\bar E_j .
      \]
      Each Yukawa term contains two matter superfields, and the \(\mu\) term
      contains none, so all have parity \(+1\).  The additional
      gauge-invariant terms are
      \[
        \mu_iL_iH_u,\quad
        \tfrac12\lambda_{ijk}L_iL_j\bar E_k,\quad
        \lambda'_{ijk}L_iQ_j\bar D_k,\quad
        \tfrac12\lambda''_{ijk}\bar U_i\bar D_j\bar D_k .
      \]
      Every one has odd matter parity.  The first three violate lepton
      number; the last violates baryon number.  Thus matter parity removes
      all renormalizable \(B\)- or \(L\)-violating superpotential terms while
      preserving the required masses and Yukawa interactions.
    \item \emph{Proton Decay from Simultaneous RPV Couplings.}
    The relevant superpotential terms may be written
      \[
        W\supset
          \lambda'_{ijk}L_iQ_j\bar D_k
          +\frac12\lambda''_{mnk}
            \epsilon^{abc}\bar U_{ma}\bar D_{nb}\bar D_{kc}.
      \]
      Exchanging \(\widetilde d_{Rk}\) produces, schematically,
      \[
        \mathcal L_{\rm eff}\supset
          \frac{\lambda'_{ijk}\lambda_{mnk}^{\prime\prime *}}
               {m_{\widetilde d_{Rk}}^2}
          \epsilon^{abc}
          \bigl(\overline{u_{Rm a}^{\,c}}d_{Rn b}\bigr)
          \left(
            \overline{\nu_{Li}^{\,c}}d_{Lj c}
            -\overline{e_{Li}^{\,c}}u_{Lj c}
          \right)+\mathrm{H.c.}
      \]
      Here charge conjugation and the chiral projections are in the
      conventional four-component notation of this volume.  The operator
      changes both \(B\) and \(L\) by one unit.  Flavor
      antisymmetry in \(\lambda''\), CKM rotations, and hadronic matrix
      elements select the physical channels, but the dangerous coefficient
      always scales as the product of couplings divided by one squark mass
      squared.
    \item \emph{Baryon Triality Check for Dimension-Four Operators.}
    Let \(\omega=e^{2\pi i/3}\), and assign additive charges modulo three
      \[
      \begin{array}{c|ccccccc}
        &Q&H_u&\bar D&L&H_d&\bar U&\bar E\\ \hline
        q&0&1&1&-1&-1&-1&-1 .
      \end{array}
      \]
      The sums for \(QH_u\bar U\), \(QH_d\bar D\),
      \(LH_d\bar E\), and \(H_uH_d\) are respectively
      \(0,0,-3,0\), so the ordinary terms are allowed.  Likewise
      \(LL\bar E\) has charge \(-3\), and \(LQ\bar D\) has charge zero.
      In contrast,
      \(\bar U\bar D\bar D\) has charge \(+1\) modulo three and is
      forbidden.  This discrete symmetry permits lepton-number violation but
      removes the renormalizable baryon-number violation responsible for the
      tree-level proton-decay catastrophe.
  \end{enumerate}

  \SupplementarySolution{2}{One-Loop Unification and Superpartner Thresholds}
  \begin{enumerate}
    \item \emph{One-Loop Supersymmetric Unification Coefficients.}
    In the convention
      \[
        \mu\frac{dg_a}{d\mu}=\frac{b_a}{16\pi^2}g_a^3,
      \]
      a Weyl fermion contributes \(2T(R)/3\), a complex scalar contributes
      \(T(R)/3\), and a gauge multiplet contributes \(-3C_2(G)\).
      Summing three generations and two Higgs doublets gives, for
      \(g_1=\sqrt{5/3}\,g'\),
      \[
        (b_1,b_2,b_3)=\left(\frac{33}{5},1,-3\right).
      \]
      Equivalently the unnormalized hypercharge coupling has \(b_{g'}=11\).
      The factor \(3/5\) follows from embedding hypercharge in a normalized
      unified generator:
      \[
        g_1^2=\frac53g'^2,\qquad
        g_3^2=g_2^2=g_1^2
      \]
      at the unification scale.  These values reproduce
      Eqs.~\eqref{eq:28.2.8}--\eqref{eq:28.2.10} with
      \(n_g=3,n_s=2\).
    \item \emph{Threshold Split Between Scalar and Fermion Partners.}
    For a representation of index \(T_i(R)\), the scalar and Weyl
      contributions to \(b_i\) are
      \[
        b_i^{s}=\frac13T_i(R),\qquad
        b_i^{f}=\frac23T_i(R).
      \]
      Since
      \(d(1/g_i^2)/d\ln\mu=-b_i/(8\pi^2)\), both particles contribute above
      \(m_f\), only the scalar contributes for \(m_s<\mu<m_f\), and neither
      contributes below \(m_s\).  Relative to placing both thresholds at
      \(m_f\), the light scalar adds
      \[
        \Delta\!\left(\frac1{g_i^2}\right)
          =\frac{T_i(R)}{24\pi^2}\ln\frac{m_f}{m_s}
      \]
      at scales below \(m_s\).  The result depends only logarithmically on
      the splitting and vanishes when the multiplet is degenerate.
  \end{enumerate}

  \SupplementarySolution{3}{Right-Handed Neutrinos and the Dimension-Five Operator}
  \begin{enumerate}
    \item In terms of \(N=N^c\), the part needed at leading order has the
      schematic form
      \[
        \mathcal L_N
          =-\frac12\bar N(\gamma^\mu\partial_\mu+M)N
           -\bar N J,
        \qquad
        J=g_\nu X+g_\nu^*X^c,
      \]
      where the second term is understood with its Majorana completion.
      The field equation gives
      \[
        N=-\frac1{\gamma^\mu\partial_\mu+M}J
          =-\frac1M J+\frac1{M^2}\gamma^\mu\partial_\mu J
           +O(M^{-3}).
      \]
      Substitution, with the factor \(1/2\) appropriate to a Majorana
      Gaussian, yields the leading effective interaction
      \[
        \mathcal L_{\rm eff}^{(5)}
          =\frac{g_\nu^2}{2M}
             X^{\mathsf T}C X+\mathrm{H.c.}
          =\frac{g_\nu^2}{2M}
            (H_i^{\mathsf T}L_j\epsilon^{ij})^{\mathsf T}
            C(H_k^{\mathsf T}L_\ell\epsilon^{k\ell})
            +\mathrm{H.c.}
      \]
      A convention that writes the Majorana bilinear without the prefactor
      \(1/2\) correspondingly assigns the coefficient \(g_\nu^2/M\); the
      physical mass below is the same.  Since
      \([H]=1\), \([L]=3/2\), this is a dimension-five operator.  The next
      term contains a derivative and is suppressed by \(M^{-2}\).

    \item At the neutral Higgs vacuum,
      \(X=-v\nu_L/\sqrt2\), up to the sign chosen for
      \(\epsilon^{ij}\).  Hence
      \[
        \mathcal L_{\rm eff}^{(5)}
          =\frac12m_\nu\nu_L^{\mathsf T}C\nu_L+\mathrm{H.c.},
        \qquad
        |m_\nu|=\frac{|g_\nu|^2v^2}{2|M|}.
      \]
      The bilinear couples \(\nu_L\) to its own charge conjugate, rather
      than to an independent right-handed light field, so the eigenstate is
      Majorana.  It carries \(\Delta L=2\), just as the heavy Majorana mass
      does.

    \item The experimental upper bound gives
      \[
        \frac{|M|}{|g_\nu|^2}
          >\frac{v^2}{2(0.5\,\mathrm{eV})}.
      \]
      Since \(0.5\,\mathrm{eV}=5.0\times10^{-10}\,\mathrm{GeV}\),
      \[
        \frac{|M|}{|g_\nu|^2}
          \gtrsim6.1\times10^{13}\,\mathrm{GeV}.
      \]
      Thus an order-one neutrino Yukawa coupling points to a heavy scale of
      at least order \(10^{14}\,\mathrm{GeV}\); a smaller \(g_\nu\) lowers
      the required \(M\) in proportion to \(|g_\nu|^2\).
  \end{enumerate}

  \SupplementarySolution{4}{Soft Flavor Tensors and Mass-Insertion Bounds}
  \begin{enumerate}
    \item \emph{Soft-Term Flavor Counting in a Super-CKM Basis.}
    After biunitary rotations diagonalize \(y_u\) and \(y_d\), with their
      left-handed mismatch placed in the CKM matrix, the soft sector still
      contains three independent Hermitian matrices
      \[
        m_Q^2,\qquad m_{\bar U}^2,\qquad m_{\bar D}^2
      \]
      and two general complex matrices \(a_u,a_d\).  Each Hermitian
      \(3\times3\) matrix contains three real diagonal entries and three
      complex off-diagonal entries; each trilinear matrix contains nine
      complex entries.  The rotations that diagonalize the Yukawa matrices
      merely transform these soft matrices by conjugation or by separate
      left and right rotations.  Gauge interactions cannot force simultaneous
      diagonalization because all three generations have identical gauge
      charges.  Alignment or degeneracy therefore requires additional
      dynamics, not standard-model gauge symmetry.
    \item \emph{A Neutral-Kaon Mass-Insertion Bound.}
    In the super-CKM basis a flavor change on a squark line is represented
      by
      \[
        (\delta^d_{12})_{LL}
          =\frac{(\Delta M_{\widetilde d}^2)_{12}}{\widetilde M^2}.
      \]
      A \(\Delta S=2\) gluino box needs one insertion on each squark
      propagator, so dimensional analysis gives
      \[
        C_{\Delta S=2}^{\widetilde g}
          \sim\frac{g_s^4}{16\pi^2\widetilde M^2}
            \bigl[(\delta^d_{12})_{LL}\bigr]^2.
      \]
      The standard-model charm box scales as
      \[
        C_{\Delta S=2}^{\rm SM}
          \sim\frac{g^4}{16\pi^2m_W^4}
            \sin^2\theta_c\cos^2\theta_c\,m_c^2.
      \]
      Requiring the first not exceed the second yields
      \[
        |(\delta^d_{12})_{LL}|
          \lesssim\frac{g^2}{g_s^2}\sin\theta_c\cos\theta_c
            \frac{m_c\widetilde M}{m_W^2},
      \]
      which is the scaling and approximate numerical bound in
      Eqs.~\eqref{eq:28.4.5}--\eqref{eq:28.4.6}.
    \item \emph{A Slepton-Mixing Bound from Radiative Muon Decay.}
    With the chirality flip placed on the external muon line, one
      off-diagonal slepton insertion gives a dipole coefficient
      \[
        C_D\sim
          \frac{eg^2}{16\pi^2}
          \frac{m_\mu}{\widetilde M^2}\,
          \delta^\ell_{12},
        \qquad
        \delta^\ell_{12}
          =\frac{(\Delta M_{\widetilde\ell}^2)_{12}}
                 {\widetilde M^2}.
      \]
      Loop functions and neutralino mixing angles are order-one modifiers.
      Since \(\Gamma(\mu\to e\gamma)\sim m_\mu^3|C_D|^2/(16\pi)\) while
      \(\Gamma_\mu\sim G_F^2m_\mu^5/(192\pi^3)\),
      \[
        {\rm Br}(\mu\to e\gamma)
          \sim\frac{\alpha g^4}{G_F^2\widetilde M^4}
            |\delta^\ell_{12}|^2
      \]
      up to numerical loop factors.  The quadratic sensitivity explains why
      even modestly light sleptons require strong degeneracy or alignment.
  \end{enumerate}

  \SupplementarySolution{5}{Physical Soft Phases and Dipole Decoupling}
  \begin{enumerate}
    \item \emph{Rephasing-Invariant Soft CP Phases.}
    Denote the Higgs bilinear coefficient by \(b=B\mu\).  Two independent
      chiral rephasings, conventionally called Peccei--Quinn and \(R\)
      rotations, can make two complex parameters real.  A convenient
      invariant basis is
      \[
        \phi_A=\arg(AM_{1/2}^*),\qquad
        \phi_B=\arg(M_{1/2}\mu b^*).
      \]
      Every field phase cancels within these products.  For example, one may
      choose \(M_{1/2}\) and \(b\) real; then \(\phi_A=\arg A\) and
      \(\phi_B=\arg\mu\).  Choosing a different convention changes the
      phases assigned to individual parameters but not these two
      combinations.  Consequently CP observables can depend on
      \(\sin\phi_A\) and \(\sin\phi_B\), never on the phase of a single
      convention-dependent parameter.
    \item \emph{Parametric Decoupling of a Supersymmetric EDM.}
    A left--right sfermion insertion has the form
      \[
        \Delta_{LR}\sim m_f(A_f-\mu^*r_f),
      \]
      with \(r_f=\cot\beta\) for up-type and \(\tan\beta\) for down-type
      fermions.  Including a gaugino mass in the chirality-changing numerator,
      the one-loop estimate is
      \[
        d_f\sim
          \frac{eg^2}{16\pi^2}
          \frac{\operatorname{Im}
            [M_\lambda\Delta_{LR}]}{\widetilde M^4}.
      \]
      Scale all dimension-one soft parameters, \(\mu\), and
      \(\widetilde M\) by \(r\).  Then
      \(M_\lambda\Delta_{LR}\) scales as \(r^2m_f\), whereas the denominator
      scales as \(r^4\).  Thus
      \(d_f\propto m_f\sin\phi/r^2\): the operator decouples quadratically,
      even though order-one phases remain.
  \end{enumerate}

  \SupplementarySolution{6}{A General Two-Higgs-Doublet Sector}
  \begin{enumerate}
    \item \emph{Charges.}
    With \(Q=T_3+Y\), the upper and lower components of a
      \(Y=-1/2\) doublet have
      \[
        Q(\phi_{\rm upper})=\frac12-\frac12=0,\qquad
        Q(\phi_{\rm lower})=-\frac12-\frac12=-1.
      \]
      Thus \(\phi=(\phi^0,\phi^-)^{\mathsf T}\).
    \item \emph{General Renormalizable Potential.}
    Define the neutral gauge singlet
      \(H\cdot\phi=H^Ti\sigma_2\phi\).  The most general potential is
      \[
      \begin{aligned}
        V={}&m_1^2H^\dagger H+m_2^2\phi^\dagger\phi
          -\left(m_{12}^2H\cdot\phi+\mathrm{H.c.}\right)\\
        &+\frac{\lambda_1}{2}(H^\dagger H)^2
          +\frac{\lambda_2}{2}(\phi^\dagger\phi)^2
          +\lambda_3(H^\dagger H)(\phi^\dagger\phi)
          +\lambda_4|H\cdot\phi|^2\\
        &+\left[
          \frac{\lambda_5}{2}(H\cdot\phi)^2
          +\lambda_6(H^\dagger H)(H\cdot\phi)
          +\lambda_7(\phi^\dagger\phi)(H\cdot\phi)
          +\mathrm{H.c.}\right].
      \end{aligned}
      \]
      Hermiticity makes
      \(m_1^2,m_2^2,\lambda_{1,2,3,4}\) real: six real parameters.
      The coefficients \(m_{12}^2,\lambda_{5,6,7}\) are complex and add
      eight more, for fourteen real coefficients in a fixed field basis.
    \item \emph{Vacuum and Gauge Spectrum.}
    Both nonzero expectation values belong to neutral components, so the
      generator \(Q=T_3+Y\) annihilates the vacuum even when \(w\ne0\).
      Electromagnetism is unbroken.  Put
      \[
        v_{\rm EW}^2=v^2+u^2+w^2.
      \]
      Expanding \(|D_\mu H|^2+|D_\mu\phi|^2\) gives
      \[
        m_W^2=\frac{g^2v_{\rm EW}^2}{4},\qquad
        m_Z^2=\frac{(g^2+g'^2)v_{\rm EW}^2}{4},\qquad
        m_\gamma^2=0.
      \]
      Of the eight real scalar components, three neutral or charged
      combinations are Goldstone modes.  The five physical modes comprise
      one charged complex scalar of charges \(+1,-1\) and three electrically
      neutral real scalars.  For generic complex coefficients and \(w\ne0\),
      the three neutral states need not have definite CP.
    \item \emph{Yukawa Couplings.}
    Let
      \(\widetilde H=i\sigma_2H^*\) and
      \(\widetilde\phi=i\sigma_2\phi^*\).  Omitting possible neutrino
      couplings, gauge invariance permits
      \[
      \begin{aligned}
        -\mathcal L_Y={}&
          \bar Q_L(Y_dH+\widehat Y_d\widetilde\phi)d_R
          +\bar Q_L(Y_u\widetilde H+\widehat Y_u\phi)u_R\\
          &+\bar L_L(Y_eH+\widehat Y_e\widetilde\phi)e_R
          +\mathrm{H.c.}
      \end{aligned}
      \]
      Each hatted and unhatted matrix is independent.  A discrete symmetry
      or alignment assumption is therefore needed if one wants to forbid
      tree-level flavor-changing neutral-scalar interactions.
  \end{enumerate}

  \SupplementarySolution{7}{Goldstone, Pseudoscalar, and Charged-Higgs Directions}
  \begin{enumerate}
    \item \emph{Neutral Goldstone and Pseudoscalar Directions.}
    Equation~\eqref{eq:28.5.26} is
      \[
        M_I^2=m_A^2
        \begin{pmatrix}
          \sin^2\beta&\sin\beta\cos\beta\\
          \sin\beta\cos\beta&\cos^2\beta
        \end{pmatrix}.
      \]
      It is the outer product of
      \((\sin\beta,\cos\beta)\) with itself.  Hence the normalized
      physical direction is
      \[
        A=\sqrt2\left(
          \sin\beta\,\operatorname{Im}\varphi_1
          \cos\beta\,\operatorname{Im}\varphi_2\right),
        \qquad m_A^2,
      \]
      while the orthogonal combination
      \[
        G^0=\sqrt2\left(
          \cos\beta\,\operatorname{Im}\varphi_1
          -\sin\beta\,\operatorname{Im}\varphi_2\right)
      \]
      has zero squared mass and is eaten by the \(Z\).  The answer follows
      from rank one, without solving a secular polynomial.
    \item \emph{The Charged-Higgs Tree-Level Sum Rule.}
    In the charge-\(-1\) basis
      \((\mathcal H_1^-,(\mathcal H_2^+)^*)\),
      Eq.~\eqref{eq:28.5.34} has the same rank-one matrix as the
      pseudoscalar sector, multiplied by \(m_A^2+m_W^2\).  Therefore
      \[
        G^-=\cos\beta\,\mathcal H_1^-
           -\sin\beta\,(\mathcal H_2^+)^*,\qquad m_{G^\pm}^2=0,
      \]
      and the orthogonal physical state
      \[
        H^-=\sin\beta\,\mathcal H_1^-
           +\cos\beta\,(\mathcal H_2^+)^*
      \]
      obeys
      \[
        m_{H^\pm}^2=m_A^2+m_W^2.
      \]
      The extra \(m_W^2\) is fixed by the \(SU(2)\) \(D\)-term quartic.
      A general two-Higgs-doublet model has independent quartics and need
      not satisfy this sum rule.
  \end{enumerate}

  \SupplementarySolution{8}{The CP-Even Higgs Sector and Its Leading Correction}
  \begin{enumerate}
    \item \emph{CP-Even Higgs Trace, Determinant, and Bounds.}
    Directly from Eq.~\eqref{eq:28.5.27},
      \[
        \operatorname{tr}M_R^2=m_A^2+m_Z^2,\qquad
        \det M_R^2=m_A^2m_Z^2\cos^22\beta.
      \]
      The quadratic formula therefore gives
      \[
        m_{H,h}^2=\frac12\left[
          m_A^2+m_Z^2
          \mathbin{\pm}
          \sqrt{(m_A^2+m_Z^2)^2
            -4m_A^2m_Z^2\cos^22\beta}
        \right].
      \]
      Since \(\sin^22\beta\geq0\), the square root is at least
      \(|m_A^2-m_Z^2|\).  The lower eigenvalue is consequently no larger
      than
      \(\tfrac12[m_A^2+m_Z^2-|m_A^2-m_Z^2|]
      =\min(m_A^2,m_Z^2)\).  Taking square roots proves the stated
      tree-level bound.
    \item \emph{Light-Higgs Alignment in the Decoupling Limit.}
    For \(m_A^2\gg m_Z^2\), write
      \[
        \alpha=\beta-\frac{\pi}{2}+\epsilon,\qquad
        \epsilon=\frac{m_Z^2}{2m_A^2}\sin4\beta
          +O(m_Z^4/m_A^4).
      \]
      The eigenvalues become
      \[
        m_h^2=m_Z^2\cos^22\beta+O(m_Z^4/m_A^2),\qquad
        m_H^2=m_A^2+m_Z^2\sin^22\beta+O(m_Z^4/m_A^2).
      \]
      Moreover \(\sin(\beta-\alpha)=1+O(\epsilon^2)\), so the light state
      points along the vacuum direction.  Its Yukawa ratios are
      \[
        \frac{g_{huu}}{g_{huu}^{\rm SM}}
          =1+\epsilon\cot\beta+O(\epsilon^2),\qquad
        \frac{g_{hdd}}{g_{hdd}^{\rm SM}}
          =1-\epsilon\tan\beta+O(\epsilon^2).
      \]
      Thus alignment is automatic as the heavy Higgs scale decouples.
    \item \emph{Leading Stop Logarithm in the Higgs Mass.}
    For degenerate stops of mass \(M_S\) and negligible left--right mixing,
      the top and stop one-loop determinants cancel above \(M_S\) but leave a
      logarithm between \(M_S\) and \(m_t\).  Differentiating the effective
      potential twice along the light Higgs direction gives
      \[
        \Delta m_h^2
          =\frac{3m_t^4}{2\pi^2v_{\rm SM}^2}
            \ln\frac{M_S^2}{m_t^2}
          =\frac{3\sqrt2G_Fm_t^4}{2\pi^2}
            \ln\frac{M_S^2}{m_t^2}.
      \]
      The color factor is three.  The correction is positive for
      \(M_S>m_t\) because supersymmetry no longer cancels the top contribution
      below the heavier scalar threshold.  Stop mixing adds important finite
      terms but does not alter this leading-log argument.
  \end{enumerate}

  \SupplementarySolution{9}{Chargino Singular Values and Limiting Spectra}
For Eq.~\eqref{eq:28.5.51},
  \[
  \begin{aligned}
    T&=\operatorname{tr}
      (\mathcal M_{\rm C}^\dagger\mathcal M_{\rm C})
      =M_2^2+|\mu|^2+2m_W^2,\\
    D&=\det(\mathcal M_{\rm C}^\dagger\mathcal M_{\rm C})
      =|M_2\mu+m_W^2\sin2\beta|^2 .
  \end{aligned}
  \]
  Hence
  \[
    m_{\chi^\pm_{1,2}}^2
      =\frac12\left[T\mp\sqrt{T^2-4D}\right],
  \]
  which expands to Eq.~\eqref{eq:28.5.52}.  If
  \(|M_2|\gg|\mu|,m_W\), the masses approach \(|M_2|\) and
  \(|\mu|\), with wino- and higgsino-like eigenstates; the roles reverse
  for \(|\mu|\gg|M_2|\).  When \(|M_2|\simeq|\mu|\), the off-diagonal
  electroweak entries cause order-\(m_W\) level repulsion and neither
  state has a pure identity.

  \SupplementarySolution{10}{Neutralino and Sfermion Mixing}
  \begin{enumerate}
    \item \emph{Neutralino Mixing in the Heavy-Higgsino Limit.}
    In the basis
      \((\widetilde B,\widetilde W^0,\widetilde H_d^0,\widetilde H_u^0)\),
      \[
        \mathcal M_N=
        \begin{pmatrix}
          M_1&0&-m_Zs_Wc_\beta&m_Zs_Ws_\beta\\
          0&M_2&m_Zc_Wc_\beta&-m_Zc_Ws_\beta\\
          -m_Zs_Wc_\beta&m_Zc_Wc_\beta&0&-\mu\\
          m_Zs_Ws_\beta&-m_Zc_Ws_\beta&-\mu&0
        \end{pmatrix}.
      \]
      Taking the Schur complement of the higgsino block gives
      \[
        \mathcal M_{\rm eff}=
        \begin{pmatrix}
          M_1-\dfrac{m_Z^2s_W^2\sin2\beta}{\mu}
          &\dfrac{m_Z^2s_Wc_W\sin2\beta}{\mu}\\[5pt]
          \dfrac{m_Z^2s_Wc_W\sin2\beta}{\mu}
          &M_2-\dfrac{m_Z^2c_W^2\sin2\beta}{\mu}
        \end{pmatrix}
        +O(\mu^{-2}).
      \]
      Away from bino--wino degeneracy, the off-diagonal term changes masses
      only at higher order, leaving the displayed diagonal shifts.
    \item \emph{Sfermion Left--Right Mixing and Fermion Mass.}
    After the Higgs fields acquire expectation values, the aligned
      trilinear term and the cross term involving \(\mu\) give
      \[
        (M_{\widetilde f}^2)_{LR}
          =m_f\left(A_f-\mu^*r_f\right),
        \qquad
        r_f=
        \begin{cases}
          \cot\beta,&f=u,c,t,\\
          \tan\beta,&f=d,s,b,e,\mu,\tau.
        \end{cases}
      \]
      Overall signs depend on the phase convention for the superpotential,
      but the factor \(m_f\) does not: both interactions carry the same
      Yukawa coupling that generates the fermion mass.  Thus first-generation
      left--right mixing is tiny even for weak-scale \(A_f,\mu\), whereas the
      stop entry can be comparable to diagonal soft squared masses.  Large
      \(\tan\beta\) can similarly enhance sbottom and stau mixing.
  \end{enumerate}

  \SupplementarySolution{11}{Messenger Stability and One-Loop Gaugino Masses}
  \begin{enumerate}
    \item \emph{Messenger Supertrace and the No-Tachyon Condition.}
    The messenger fermions form a Dirac pair with
      \[
        m_\psi=|\lambda\mathcal S|.
      \]
      The scalar quadratic form is diagonalized by equal mixtures of
      \(\phi\) and a phase-rotated \(\bar\phi^*\), giving
      \[
        m_\pm^2=|\lambda\mathcal S|^2
          \pm|\lambda\mathcal F|.
      \]
      Counting real degrees of freedom,
      \[
        \operatorname{STr}M^2
          =2(m_+^2+m_-^2)-4m_\psi^2=0.
      \]
      The lighter scalar is non-tachyonic only if
      \[
        |\lambda\mathcal S|^2\geq|\lambda\mathcal F|
        \quad\Longleftrightarrow\quad
        |\mathcal F|\leq|\lambda|\,|\mathcal S|^2.
      \]
      Equality produces a massless messenger scalar; a controllable
      messenger expansion normally requires a strict inequality.
    \item \emph{Nonidentical Spurions and One-Loop Gaugino Masses.}
    The holomorphic threshold of messenger pair \(n\) is
      \(\ln(\lambda_nS_n)\).  Its \(\mathcal F\)-component is
      \(\mathcal F_n/\mathcal S_n\), so for gauge factor \(a\)
      \[
        M_a=\frac{g_a^2}{16\pi^2}
          \sum_n2T_a(R_n)\frac{\mathcal F_n}{\mathcal S_n}.
      \]
      Here \(T_a(R_n)\) is the Dynkin index of one member of the conjugate
      pair; for grand-unified \(U(1)\),
      \(T_1=(3/5)\sum_{\rm states}Y^2\).
      The sum is complex, so different spurion phases can interfere and the
      physical mass is \(|M_a|\).  The Yukawa coefficients \(\lambda_n\)
      cancel at leading order.  Corrections begin at second order in
      \(|\mathcal F_n|/(|\lambda_n||\mathcal S_n|^2)\).
  \end{enumerate}

  \SupplementarySolution{12}{Scalar Soft Masses and Unified Messenger Indices}
  \begin{enumerate}
    \item \emph{Two-Loop Scalar Masses from Quadratic Casimirs.}
    In minimal gauge mediation,
      \[
        m_{\widetilde f}^2
          =2\sum_a C_a(f)
           \left(\frac{g_a^2}{16\pi^2}\right)^2
           \sum_n2T_a(R_n)
           \left|\frac{\mathcal F_n}{\mathcal S_n}\right|^2.
      \]
      The relevant Casimirs in \(SU(5)\)-normalized hypercharge are
      \[
      \begin{array}{c|ccc}
        &C_3&C_2&C_1\\ \hline
        Q&4/3&3/4&1/60\\
        \bar U&4/3&0&4/15\\
        L&0&3/4&3/20\\
        \bar E&0&0&3/5 .
      \end{array}
      \]
      Substitution yields the corresponding masses.  The formula predicts
      heavy colored scalars, intermediate weak doublets, and the lightest
      gauge contribution for the hypercharge-only singlet.  Positivity and
      generation independence follow directly from the absolute squares and
      gauge Casimirs.
    \item \emph{Complete Messenger Multiplets and Unification.}
    The decomposition is
      \[
        \boldsymbol5=(\boldsymbol3,\boldsymbol1,-1/3)
          \oplus(\boldsymbol1,\boldsymbol2,+1/2),
      \]
      together with the conjugate representations.  A representation plus
      its conjugate has
      \[
        n_3=2T(\boldsymbol3)=1,\qquad
        n_2=2T(\boldsymbol2)=1.
      \]
      For hypercharge, summing over every component gives
      \[
        n_1=\frac35\left[
          2(3)\left(\frac13\right)^2
          +2(2)\left(\frac12\right)^2\right]
          =1.
      \]
      Thus \(\Delta b_1=\Delta b_2=\Delta b_3=1\).  Above a common messenger
      threshold all three inverse couplings acquire the same extra
      logarithmic slope in unified normalization, so their one-loop meeting
      point is preserved, although the common unified coupling changes.
  \end{enumerate}

  \SupplementarySolution{13}{Flavor Universality and the Gravitino LSP}
  \begin{enumerate}
    \item \emph{Flavor Universality at the Messenger Scale.}
    Gauge vertices depend on representation matrices but not on generation.
      Therefore the two-loop result has the form
      \[
        (m_{\widetilde f}^2)_{ij}
          =\delta_{ij}\,2\sum_aC_a(f)
           \left(\frac{g_a^2}{16\pi^2}\right)^2\Lambda_a^2
      \]
      at the messenger scale.  RG evolution introduces Yukawa matrices,
      especially \(y_t\), through combinations such as
      \(y_uy_u^\dagger m_Q^2\) and
      \(y_u m_{\bar U}^2y_u^\dagger\).  These split the third generation
      from the first two and induce CKM-suppressed off-diagonal terms.
      Crucially, the new flavor tensors are built from the same Yukawas that
      break flavor in the standard model.  Gauge mediation therefore begins
      universal and its leading departures have minimal-flavor-violation
      alignment.
    \item \emph{Gravitino Scaling and the Identity of the LSP.}
    If the same order parameter \(F\) controls both sectors,
      \[
        m_{3/2}\sim\frac{F}{\sqrt3M_{\rm Pl}},\qquad
        m_{\rm soft}\sim\frac{\alpha}{4\pi}\frac{F}{S}.
      \]
      Their ratio is
      \[
        \frac{m_{3/2}}{m_{\rm soft}}
          \sim\frac{4\pi}{\sqrt3\alpha}\frac{S}{M_{\rm Pl}}.
      \]
      Hence a messenger scale well below
      \((\alpha/4\pi)M_{\rm Pl}\) makes the gravitino the LSP.  Goldstino
      equivalence gives an NLSP width of order
      \[
        \Gamma_{\rm NLSP}\sim\frac{m_{\rm NLSP}^5}{16\pi F^2},
        \qquad
        c\tau\sim\frac{16\pi F^2}{m_{\rm NLSP}^5}.
      \]
      The lifetime is therefore exceptionally sensitive to the
      supersymmetry-breaking scale even when the visible spectrum is fixed.
  \end{enumerate}

  \SupplementarySolution{14}{The Gauge-Mediation \(\mu\)--\(B\mu\) Problem}
If a product of messenger couplings \(c\) generates the Higgs bilinear
  at one loop, dimensional analysis gives
  \[
    \mu\sim\frac{c}{16\pi^2}\frac{F}{S},\qquad
    B\mu\sim\frac{c}{16\pi^2}\left(\frac{F}{S}\right)^2.
  \]
  Therefore
  \[
    \frac{B\mu}{\mu^2}\sim\frac{16\pi^2}{c},
  \]
  which is parametrically much larger than unity for perturbative
  \(c\).  Electroweak minimization instead needs \(B\mu\) comparable to
  the Higgs soft squared masses and \(\mu^2\), all of weak-scale size.
  The same mechanism that produces a viable \(\mu\) thus typically
  overproduces \(B\mu\) by a loop factor.  Separating their origins or
  adding symmetries and additional dynamics is necessary.

  \SupplementarySolution{15}{Unitary Gauge for a General Higgs Mechanism}
  \begin{enumerate}
    \item Put \(g=g_\phi e^{i\epsilon^aT_a}\).  At the maximum,
      \[
        0=\left.\frac{\partial f_\phi}{\partial\epsilon^a}\right|_0
          =-\operatorname{Im}(\widetilde\phi^\dagger\tau_av),
      \]
      where unitarity was used to move \(D(g_\phi)\) from \(v\) to
      \(\phi\).  For \(\phi=v\), the identity is a maximum by
      Cauchy--Schwarz, so one may take \(g_\phi=1\).  Gauge invariance then
      implies \(\langle\widetilde\phi\rangle=v/\sqrt2\).

    \item Since \(D(h)v=v\),
      \(f_\phi(g_\phi h)=f_\phi(g_\phi)\) for every \(h\in H\).  The
      ambiguity is therefore exactly a residual local \(H\) transformation.
      If \(\widetilde A_\mu=A_\mu^{\mathfrak h}+A_\mu^\perp\), then
      \[
        A_\mu^{\mathfrak h}\longmapsto
          h^{-1}A_\mu^{\mathfrak h}h+i h^{-1}\partial_\mu h,
        \qquad
        A_\mu^\perp\longmapsto h^{-1}A_\mu^\perp h.
      \]
      Thus the unbroken components remain connections, while the broken
      components form a linear representation of \(H\).

    \item The expansion is
      \[
      \begin{split}
        f_{\widetilde\phi}(e^{i\theta\cdot T})
          ={}&\operatorname{Re}(\widetilde\phi^\dagger v)
          -\theta^a\operatorname{Im}
             (\widetilde\phi^\dagger\tau_av)\\
          &-\frac12\theta^a\theta^b
             \operatorname{Re}
             (\widetilde\phi^\dagger\tau_a\tau_bv)+O(\theta^3).
      \end{split}
      \]
      The linear term vanishes by part (a).  At
      \(\widetilde\phi=v/\sqrt2\), the quadratic change is
      \[
        -\frac1{2\sqrt2}\left\|\theta^a\tau_av\right\|^2.
      \]
      Continuity preserves strict negativity near the vacuum in every
      direction transverse to \(\mathfrak h\).  Its only zero directions
      obey \(\theta^a\tau_av=0\), hence lie in \(\mathfrak h\).  The
      representative is consequently locally unique on \(G/H\).

    \item Substitution gives
      \[
      \begin{split}
        \mathcal L={}&-\frac14\sum_a\frac1{g_a^2}
          \widetilde F^a_{\mu\nu}\widetilde F_a^{\mu\nu}
          -\frac12(\partial_\mu h)^\dagger\partial^\mu h
          -V\!\left(\frac{v+h}{\sqrt2}\right)\\
        &+\frac i2\widetilde A^a_\mu
          \left[(\partial^\mu h)^\dagger\tau_a(v+h)
               -(v+h)^\dagger\tau_a\partial^\mu h\right]\\
        &-\frac12\widetilde A^a_\mu\widetilde A^{b\mu}
             (v+h)^\dagger\tau_a\tau_b(v+h).
      \end{split}
      \]
      Only the symmetric part in \(a,b\) contributes to the last line.
      With \(\widetilde A^a_\mu=g_a\widehat A^a_\mu\), its vacuum term is
      \(-\tfrac12\widehat A^a_\mu(\mu^2)_{ab}\widehat A^{b\mu}\), where
      \[
        (\mu^2)_{ab}=\frac{g_ag_b}{2}
        v^\dagger\{\tau_a,\tau_b\}v.
      \]
      Hermiticity makes this real and symmetric, and
      \[
        c^a(\mu^2)_{ab}c^b
          =\left\|\sum_a g_ac^a\tau_av\right\|^2\geq0.
      \]
      It vanishes exactly when the corresponding generator annihilates
      \(v\), which is precisely the Lie algebra of \(H\).

    \item Write \(\phi=\rho e^{i\vartheta}\).  For charge \(p\), one may
      choose \(g_\phi=e^{i\vartheta/p}\), modulo the residual subgroup
      \(\mathbb Z_{|p|}\), and obtains \(\widetilde\phi=\rho\).  Since
      \(\tau=p\),
      \[
        \mu^2=g^2p^2v^2.
      \]
      This is the usual abelian unitary gauge: the phase is removed and its
      degree of freedom becomes the longitudinal polarization of the vector.

    \item Take \(\tau_i=\sigma_i/2\) and
      \(\tau_Y=-\mathbf1/2\).  The maximizing condition removes the three
      Goldstone directions and leaves
      \[
        \widetilde\phi=\frac1{\sqrt2}\binom{v+h}{0}.
      \]
      Inserting the generators into the general formula gives
      \[
        \mu^2=\frac{v^2}{4}
        \begin{pmatrix}
          g_2^2&0&0&0\\
          0&g_2^2&0&0\\
          0&0&g_2^2&-g_1g_2\\
          0&0&-g_1g_2&g_1^2
        \end{pmatrix}.
      \]
      Thus \(W^\pm=(\widehat A^1\mp i\widehat A^2)/\sqrt2\) have
      \(m_W^2=g_2^2v^2/4\).  The neutral combinations are
      \[
        Z_\mu=\frac{g_2\widehat A^3_\mu-g_1\widehat B_\mu}
                    {\sqrt{g_1^2+g_2^2}},
        \qquad
        A_\mu=\frac{g_1\widehat A^3_\mu+g_2\widehat B_\mu}
                    {\sqrt{g_1^2+g_2^2}},
      \]
      with \(m_Z^2=(g_1^2+g_2^2)v^2/4\) and \(m_A=0\).
  \end{enumerate}

\end{enumerate}
